\section{进入复域}

微积分学在建立了以后，立即就转到复域中了，这个发展是很自然的，又是很重要的。这样也就形成了数学的一个广阔的领域——复分析。特别是在 19 世纪，几乎所有重要的数学家都对它有贡献。许多人都有一个看法，即复分析是当时数学的中心，其所以如此，一方面来自它对解决实际的物理、工程问题的效用，另一方面则是因为许多在实数域中十分复杂困难的问题在复域中变简单了，其本质更出现了。因此复分析就成了一个很完美的框架。在本书里，我们总是随时准备进入复域，首先讲某一个问题中如何直接应用实域中的微积分理论，一直到遇到了本质上与实域不同之处才止步。这样为读者将来学习复分析作一个准备。

于是关于函数，读者都知道，狄利克雷关于函数的定义是：设有两个变量 $x$ 与 $y$，$x$ 在某一集合 $A$ 中变化，而对每个 $x\in A$，均有一个 $y$ 与之相对应，则说 $y$ 是 $x$ 的一个函数（准确些说是 $x$ 的定义于 $A$ 上的函数），记作 $y=f(x)$，而 $A$ 称为 $f$ 的定义域。这个定义中“一个”两字十分重要，我们下面要仔细讨论。除此以外，有两点值得注意，首先是 $x$ 与 $y$ 也可是复数（但他没有明说）。这样不会造成什么困难，我们只需规定 $x$ 与 $y$ 都是复数即可。然而进一步的研究表明，$x$ 是否为“真正的复数”（即其虚部是否真正为 0），有极大的关系。以下，凡讲到复变量的函数时，总认为 $x$ 是真正的复数，而不止是作为复数的特例的实数。这时应该更准确地说是复自变量函数。至于 $y$ 是否真正的复数并无关系，即令 $x$ 为实而 $y$ 为复，我们将得到一个“取复值的函数”。这时不妨分开 $y$ 的实、虚部，并且研究两个取实值的函数。反过来 $x$ 是否真正的复数大有关系，而 $y$ 是否真正的复数则没有关系（只不过在下面讨论的所谓解析函数，当 $x$ 可以取复值时，$y$ 不可能只取实值，除非 $y=$ 常数。而在复域中我们又只讨论解析函数）。以后，凡考虑复变量函数时，自变量与函数值总用 $z$ 和 $w$ 记：$w=f(z)$。

另一点值得注意的是，在狄利克雷的定义中 $A$ 是任意集合，但是狄利克雷当时考虑的实际问题中 $A$ 则是一个区间，所以原定义中是说 $x$ 在一区间中变动。就在那时，人们开始认识，有必要考虑任意集合上定义的函数，所以现在的微积分学中讲函数定义时，定义域总是一般的集合。对复变函数自然也是一样。但是我们要讨论的复变函数都是解析函数，需要对它进行微分与积分，这时，需要假设 $A$ 为开集，因为要避免在 $A$ 的边界上讨论这些问题。我们常说 $A$ 是一个区域。区域就是连通开集，但有的书上用语不太注意时把开集也就说成了区域。以后在讨论复变量的函数时，我们总规定 $A$ 是区域，整个复平面是一个区域，复平面除去有限多个点也是区域，当然还有更复杂的区域。

有关函数和变量的一些基本概念如极限、连续性、无穷小量以及连续函数的一些基本性质，实的情况与复的情况是一样的。但是因为复数不能比较大小，所以涉及比较大小的概念如单调性、上下确界、最大最小值等等都没有意义。但是 $|f(z)|$ 却是 $z=x+iy$ 的实值函数，也就是 $x,y$ 的实值函数，所以若 $f(z)$ 在有界闭区域 $G$ 上连续，则 $|f(z)|$ 必可在 $G$ 上达到最大最小值，可以取中间值等等，但是对 $f(z)$ 就谈不上这些问题。在涉及微分积分时，如前面所指出的，只需注意复常数可以如实常数一样看待，从微分与积分号下提出来或放进去，均与实的情况一样。当然，中值定理因涉及比较大小就不可能成立了。这些问题在下面两章讨论复变量函数的微分与积分时都会详细讨论，这里就一句带过。

我们最要注意的是“多值函数”的问题。按函数的定义，所谓 $y=f(x)$ 是 $x$ 在定义域 $A$ 上的函数，即指对于一个 $x\in A$，必有“一个”$y$ 与之相应。上文我们特别指出“一个”两字的重要性，所以按函数的定义，函数就是单值函数，无所谓多值函数。

这样的限制是有道理的。如果我们允许对同一个 $x$ 有多个 $y$ 与之相应，则下面的“东西”算不算一个多值函数？
\begin{equation}
 y=\begin{cases}
 \sin x,\\
 \tan x,\\
 x^3+e^x,\\
 \cdots
 \end{cases}
 \tag{1}
\end{equation}
这样把一些毫不相干的东西硬拉在一起并称之为“多值函数”的作法，至少是没有益处的。

但是在数学中给出一个定义就是划了一条界线，在界线内的对象是需要研究的，而在界线外的如上面的例子就可能是需要排除的。因此界线的划法就很有讲究，应该要能保证数学研究的进展。现在我们要问，在给出了函数的定义如上以后，我们在数学的研究与应用中常用的多值函数例如 $w=\sqrt[n]{z}$ 应如何对待？

我们不要把函数按定义必须是单值函数这一点看死了。在现代数学中有所谓集值映射（set-valued map）的概念，它在例如经济数学中就很有用。它的定义是，对每个 $x\in A$，必有某个“大集合”即一个空间 $Y$ 的子集 $y$（注意现在是 $y\subset Y$，而不是 $y\in Y$）与之对应，$y$ 中可能有许多元素如 $y=\{y_1,y_2,\cdots\}$（甚至不是可数多个），$y_i\in Y$ 而不是 $y_i\subset Y$。这个概念的来源之一就是如同 $w=\sqrt[n]{z}$ 这样的“多值函数”。集值映射这个概念常见于“非线性分析”这个很新很活跃的分支，我们不能涉及，只说一下，读完本书对进入这个领域大有好处。

我们现在讨论 $w=\sqrt[n]{z}$ 却是从一个很古老的观点出发的：把它看成 $w=z^n$ 的反函数。这里的记号要解释一下，$w=\sqrt[n]{z}$ 是 $z=w^n$ 的反函数，为什么又说 $w=\sqrt[n]{z}$ 又是 $w=z^n$ 的反函数呢？要注意重要的是“函数关系”：“$n$ 次根”函数关系与“$n$ 次幂”函数关系互为反函数。如果某一个函数关系记为 $f$，而且已知有反函数存在，则反函数关系，亦即逆映射记为 $f^{-1}$。于是 $f\circ f^{-1}=id$，$f^{-1}\circ f=id$。$id$ 是恒等映射。不过这两个 $id$ 并不相同，因为其定义域不同。前一个 $id$ 定义在 $f$ 的值域上，后一个定义在 $f$ 的定义域上。总之我们现在讨论 $w=\sqrt[n]{z}$。若固定 $z$ 值，它是 $n$ 次代数方程
\begin{equation}
 w^n-z=0
 \tag{2}
\end{equation}
的根。这是一个很古老的问题，当年高斯第一次证明了 $n$ 次代数方程\newline $\sum_{k=0}^{n}a_{n-k}w^k=0\ (a_0\ne0)$ 必有 $n$ 个复根存在，这是数学史上第一次有了一个“纯粹的存在定理”。不过我们现在要讨论的是更一般的方程
\begin{equation}
 a_0(z)w^n+a_1(z)w^{n-1}+\cdots+a_n(z)=0,
 \tag{3}
\end{equation}
而系数 $a_i(z)$ 都是 $z$ 的多项式。从这种方程解出的 $w=w(z)$ 称为代数函数，它们大多是多值的，这个问题在 19 世纪中后期是数学中的重大问题，黎曼正是由研究它们才得到十分深刻的黎曼曲面的思想。所谓黎曼曲面 $S$ 是这样一种“曲面”，使得上述方程的“多值”的解成为 $S$ 上的单值函数。由于(2)是(3)的特例，我们现在就以(2)为特例介绍这个极为深刻的思想。

对于 $z=1$，方程(2)有 $n$ 个根\jiaokan{原文作“对于固定的 $z$”，但紧接的式(4)列出的是 $w^n=1$ 的 $n$ 个单位根；一般固定 $z$ 的 $n$ 个根在后面的式(5)给出。}
\begin{equation}
 w^{(k)}=e^{i2k\pi/n},\qquad k=0,1,\cdots,n-1,
 \tag{4}
\end{equation}
如果让 $k$ 取其他整数值如 $k=n,n+1,\cdots$ 或 $k=-1,-2,\cdots$，所得的 $w^{(k)}$ 仍然在(4)式所给的 $n$ 个 $w^{(k)}$ 之中，(4)中的 $n$ 个 $w^{(k)}$ 称为 1 的 $n$ 个 $n$ 次单位根，它们位于一个正 $n$ 边形的 $n$ 个顶点上（如图2-3-1），其中一个是 1（即 $w^{(0)}$）。

\begin{figure}[htbp]
\centering
\begin{tikzpicture}[bella figure,scale=1.0,>=stealth]
  \def\R{1.65}
  \draw (0,0) circle (\R);
  \foreach \k in {0,...,5}{\coordinate (P\k) at ({60*\k}:\R);}
  \draw[thick] (P0)--(P1)--(P2)--(P3)--(P4)--(P5)--cycle;
  \foreach \k in {0,...,5}{\draw (0,0)--(P\k);}
  \fill (0,0) circle (1.2pt); \node[below left=1pt] at (0,0) {$O$};
  \node[right=4pt] at (P0) {$e^{i2\pi}=1$};
  \node[above right=3pt] at (P1) {$e^{i2\pi/n}$};
  \node[below right=3pt] at (P5) {$e^{i2\pi(n-1)/n}$};
  \node[right=2pt] at (0.75,0.18) {$2\pi/n$};
\end{tikzpicture}
\caption*{图2-3-1}
\end{figure}

我们可以再叙述一次上面的结果。由欧拉公式，任何复数 $z$ 均可写为 $z=re^{i\varphi}$，这里 $r\ge0$，$\varphi$ 称为辐角，它是多值的，因为若 $\varphi$ 是一个辐角，则 $\varphi+2k\pi,k=0,\pm1,\pm2,\cdots$ 都是辐角。但 $r=0$ 时辐角不确定。容易证明，(2)的 $n$ 个根即
\begin{equation}
 w^{(k)}=z^{1/n}=r^{1/n}e^{i(\varphi+2k\pi)/n},\qquad k=0,1,\cdots,n-1.
 \tag{5}
\end{equation}

$r^{1/n}$，也就是正实数 $r$ 的 $n$ 次算术根。如果 $z$ 是一正实数，则可取辐角 $\varphi=0$，这时在上述 $n$ 个根中有一个是正实数（相应于 $k=0$），这个根也就称为 $z^{1/n}$ 的算术根。而对于一般的复数，算术根的概念是没有意义的。有的书上规定只有算术根才可使用 $\sqrt[n]{z}$ 的记号，一般情况下则使用 $z^{1/n}$，我们却不如此拘谨，主要是要知道它是“多值函数”，而且这种多值性的来源就在于复数辐角的多值性。为了从这样多值中分出特别的一个便于使用，我们时常指定辐角 $\varphi$ 的取值范围是 $0\le\varphi<2\pi$，并称这个值为“主值”。但是这又是一个不必过分计较的名词，有时也说 $-\pi\le\varphi<\pi$ 是主值，$-\pi<\varphi\le\pi$ 是主值，$-\pi\le\varphi\le\pi$ 是主值……总之，目的在于分离出一个特定的值来。

乍看一下，$w^n=z$ 有 $n$ 个反函数 $w^{(k)}$，但仔细看一下又不能泰然处之，因为若令 $z_0=r_0e^{i\varphi_0}$，$0\le\varphi_0<2\pi$，并取 $w^{(0)}(z_0)=r_0^{1/n}e^{i\varphi_0/n}$ 为 $w^{(0)}(z)$ 之“初始值”（$w^{(0)}(z_0)$ 不妨称为 $w=z^{1/n}$ 的主值），现在来看 $z$ 变化时 $w$ 之值如何变化。这时 $z$ 变动的道路(path)之构造将起关键作用。例如取图2-3-2上的不绕过 $O$ 的道路 $L_1$，如果 $z$ 从 $z_0$ 开始依逆时针方向绕 $L_1$ 一圈回到 $z_0$，则 $z$ 变回原来的 $r_0e^{i\varphi_0}$，而 $w$ 也由 $w^{(0)}(z_0)$ 变回原来的 $w^{(0)}(z_0)$。但是 $z$ 从 $z_0$ 开始沿绕过 $O$ 的道路 $L_2$ 依逆时针方向又回到 $z_0$，后一个 $z_0$ 将是 $r_0e^{i(\varphi_0+2\pi)}$，而 $w$ 将由 $w^{(0)}(z_0)$ 变成 $r_0^{1/n}e^{i(\varphi_0+2\pi)/n}$，即变为 $w^{(1)}(z_0)$。若 $z$ 再沿 $L_2$ 绕 $O$ 一圈再次回到 $z_0$，现在的 $z_0$ 将是 $r_0e^{i(\varphi_0+2\pi\cdot2)}$，而 $w$ 由 $w^{(1)}(z_0)$ 变成 $w^{(2)}(z_0)=r_0^{1/n}e^{i(\varphi_0+2\cdot2\pi)/n}$。所以这 $n$ 个 $w=z^{1/n}$ 将依次互相转变，所以把 $w=z^{1/n}$ 看成一个单值的反函数更为妥当，但是现在给定一个 $z$，可以求出 $n$ 个 $w$ 值如(5)，这个矛盾如何解决？

\begin{figure}[htbp]
\centering
\begin{tikzpicture}[bella figure,scale=0.95,>=stealth]
  \draw[->] (-2.1,0)--(2.35,0) node[right=2pt] {Re $z$};
  \draw[->] (0,-1.85)--(0,2.15) node[above=2pt] {Im $z$};
  \node[below left=1pt] at (0,0) {$O$};
  \coordinate (Z) at (1.45,1.35);
  \fill (Z) circle (1.5pt); \node[above right=2pt] at (Z) {$z_0$};
  \draw[thick,->] (Z) .. controls (1.9,0.8) and (1.35,0.1) .. (0.75,0.35)
        .. controls (0.15,0.65) and (0.45,1.45) .. (Z);
  \node[right=3pt] at (0.82,0.52) {$L_1$};
  \draw[thick,->] (Z) .. controls (0.45,2.0) and (-1.65,1.35) .. (-1.7,0)
        .. controls (-1.65,-1.2) and (-0.5,-1.65) .. (0.65,-1.3)
        .. controls (1.45,-1.0) and (2.0,0.5) .. (Z);
  \node[above left=2pt] at (-1.35,0.95) {$L_2$};
\end{tikzpicture}
\caption*{图2-3-2}
\end{figure}

办法就是把原来的 $z_0$ 与绕 $L_2$ 转了一圈使 $\varphi_0$ 增加了 $2\pi$ 的 $z_0=r_0e^{i(\varphi_0+2\pi)}$ 区别开来。这里就表现了黎曼的天才：设想有 $n$ 个 $z$ 平面，如同 $n$ 张纸一样，而且把它们钉在一起，可是不是像书那样沿书脊钉，而且在 $z=0$ 处钉住，然后例如把它们都沿正实轴剪开，再把一张纸的沿正实轴的下沿（辐角接近 $2\pi$ 处）与另一张纸的上沿（辐角接近 0 处）粘在一起，像一个旋转楼梯。当我们沿楼梯旋转一周环绕 $O$ 点就自动上了一层楼。初始的 $z_0$ 在一楼，下一个 $z_0$ 在二楼……这样就把两个 $z_0$ 区别开来了。这样一层又一层地转，$w^{(0)}$ 就自动地变成 $w^{(1)}\cdots$ 再变成 $w^{(n)}$ 等等。可是到了 $n$ 楼（即 $k=n-1$ 处），如果再转，$w^{(n-1)}$ 就应变成 $r_0^{1/n}e^{i(\varphi_0+2(n-1)\pi+2\pi)/n}=r_0^{1/n}e^{i\varphi_0/n}$，所以 $z_0$ 就应该回到一楼的起点。图2-3-3，每层 $z_0$ 怎么能由 $n$ 楼神不知鬼不觉地再回起点呢？$z_0$ 的路径应该与其它楼面都相交，可是交点处的 $z_0$ 对应的 $w$ 值又怎么计算呢？可是黎曼认为不必管，就让 $z$ 像游山逛水那样钻天入地了无痕迹算了，黎曼把这样造出来的东西 $S_n$ 看成一个曲面，后世就称为黎曼曲面。于是所谓多值函数 $w=z^{1/n}$ 就成了黎曼曲面 $S_n$ 上的单值函数。自此以后，研究多值函数就变成了研究如何构造相应的黎曼曲面。

\begin{figure}[htbp]
\centering
\begin{tikzpicture}[bella figure,x=1cm,y=0.68cm,>=stealth]
  \foreach \k/\lab in {0/{k=0},1/{k=1},2/{},3/{},4/{k=n-1}}{
    \draw[thick] (-2.0,\k)--(-0.55,\k) .. controls (-0.1,\k) and (0.2,{\k+0.78}) .. (0.62,{\k+0.78}) -- (2.25,{\k+0.78});
    \ifx\lab\empty\else\node[right=3pt] at (2.25,{\k+0.78}) {$\lab$};\fi
  }
  \draw[dashed] (0.05,-0.35)--(0.05,5.15);
  \node[below=3pt] at (0.05,-0.35) {$z=0$};
  \node[right=4pt] at (2.25,-0.15) {$z$ 平面};
\end{tikzpicture}
\caption*{图2-3-3}
\end{figure}

当然，读者会感到困惑，$S_n$ 究竟是怎样构造的呢？我们不妨这样来作。先找一个小纸片，在上面标记一点 $z_0$，然后例如沿着 $L_2$ 与 $z_0$ 相近处，再找一个小纸片与第一片粘起来，粘了一整圈以后，当 $L_2$ 回到 $z_0$ 时，注意在这时的小纸片上虽然仍然标记上 $z_0$，可是不要与原来最初的小纸片粘在一起，而要用新的小纸片粘下去，这样到了第 $n$ 圈时，已经用过了不少小纸片，后来的小纸片与第二片、第三片大相径庭，所以可以“置之不理”，而硬是把它与最初的一片粘在一起。这有点像运动会上的万米赛跑，运动员要在 400 米的圆形跑道上跑 25 圈。于是当一个运动员一路领先已到了冲刺阶段，金牌在望时，突然看见前面还有一个运动员，他当然可以“置之不理”，因为“前面”的运动员还只跑了 24 圈，还差 400 米呢。上面我们写了一些“游戏文章”，其实例如道路、粘结、旋转楼梯等等都有数学文章在后面。

\begin{figure}[htbp]
\centering
\begin{tikzpicture}[bella figure,scale=0.92,>=stealth]
  \coordinate (O) at (0,0);
  \draw[thick] (O)--(3.2,0) node[right=3pt] {$A_0$};
  \draw[thick] (O)--(2.75,1.55) node[right=3pt] {$A_1$};
  \draw[thick] (O)--(1.55,2.75);
  \draw[dashed] (O)--(2.6,-1.0);
  \node[left=2pt] at (O) {$O$};
  \node at (1.2,0.28) {I}; \node at (1.02,0.88) {II}; \node at (0.55,1.55) {III}; \node at (1.0,-0.35) {IV};
  \node[below=2pt] at (1.65,0) {$N$};
  \draw[thick] (1.8,-0.45) .. controls (2.85,-0.3) and (3.05,0.8) .. (2.35,1.15)
    .. controls (1.45,1.55) and (1.45,2.7) .. (2.15,3.1);
  \draw[thick] (1.25,0.25) .. controls (2.4,0.25) and (2.5,1.0) .. (2.1,1.35)
    .. controls (1.7,1.75) and (2.05,2.3) .. (2.8,2.35);
  \draw[thick] (2.1,0.18) .. controls (3.0,0.1) and (3.0,0.85) .. (2.45,1.0);
  \fill (2.45,0.92) circle (1.4pt);
\end{tikzpicture}
\caption*{图2-3-4}
\end{figure}

如果读者仍然感到 $S_n$ 神秘莫测，我们再换一个角度来看：$w=z^{1/n}$ 是 $w=z^n$ 的反函数，因此其定义域是后者的值域。$w=z^n$ 映 $\mathbb C$（$z$ 平面）$\to S_n$，我们现在研究 $S_n$ 的构造。把 $z$ 平面 $C$ 分成 $n$ 个扇形，其顶角均为 $2\pi/n$，于是 $\overline{OA_0}$ 被映为 $w$ 平面的正实轴，$\overline{OA_1}$ 也是一样，总之第一个扇形被一对一地映为 $w$ 的全平面，而 $\overline{OA_0}$ 我们认为变为 $w$ 平面正实轴“上岸”，$\overline{OA_1}$ 变为 $w$ 平面正实轴“下岸”，总之，我们得到第一页的平面。再往下进行到了第二个扇形，$\overline{OA_1}$ 变成第二页 $w$ 平面正实轴的上岸，而与第一页相邻。仿此进行下去可以作出 $n$ 页 $w$ 平面而且再回到第一页。在图2-3-4上我们画了一些小区域，它们都是跨在两页 $w$ 平面上的，借此标志第一页与第二页以及第 $n$ 页相邻，但与第三页、第 $n-1$ 页都不相邻。所以在第 $n$ 页上的 $w$ 又会“从天而降”，不知怎么又回到了第一页。这样看来 $w=z^{1/n}$ 的值域 $S_n$ 显然具有上面讲的黎曼曲面的构造，所以它的反函数 $w=z^{1/n}$ 应以这样的 $S_n$ 为定义域，而且成为 $S_n$ 上的单值函数。所以，黎曼曲面是实实在在的结构而非空中楼阁！

有了对 $w=\sqrt[n]{z}$ 的详细讨论，再看 $w=e^z$ 的反函数 $w=\ln z$ 就容易多了。事实上，因为
\[
z=re^{i\varphi}=re^{i(\varphi+2k\pi)},\qquad k=0,\pm1,\pm2,\cdots,
\]
\begin{equation}
 w=\ln z=\ln r+i(\varphi+2k\pi),\qquad k=0,\pm1,\pm2,\cdots.
 \tag{6}
\end{equation}
从表面上看来，$\ln z$ 是一个无穷值的多值函数，但是如果按上面讨论 $w=\sqrt[n]{z}$ 的方法来处理，可以知道它只是一个定义在黎曼曲面 $S_\infty$ 上的函数。$S_\infty$ 由无穷多页 $z$ 平面叠合而成（图2-3-5）。只不过，它一方面可以向上（如果 $z$ 绕 $z_0$ 作反时针方向旋转的话）升到无穷多层，也可以向下（如果 $z$ 绕 $z_0$ 作顺时针方向旋转的话），降到无穷多层的“地下室”去。读者由此可见，$w=\ln z$ 是黎曼曲面 $S_\infty$（它有无穷多页）上的单值函数。

\begin{figure}[htbp]
\centering
\begin{tikzpicture}[bella figure,x=1cm,y=0.60cm,>=stealth]
  \foreach \k in {-2,-1,0,1,2}{
    \draw[thick] (-2.15,\k)--(-0.55,\k) .. controls (-0.15,\k) and (0.15,{\k+0.62}) .. (0.65,{\k+0.62}) -- (2.25,{\k+0.62});
  }
  \draw[dashed] (0.05,-2.55)--(0.05,3.15);
  \node[right=3pt] at (2.25,0.62) {$z=0$};
\end{tikzpicture}
\caption*{图2-3-5}
\end{figure}

还有反三角函数也是相应黎曼曲面上的单值函数，但因其黎曼曲面的构造比较复杂，我们就不再讨论了。

可是无论怎么说，神秘气氛终于难以尽除。为什么数学家们肯接受这样的东西？根本之点在于它能解决深刻的数学问题和物理问题。从上面讲的处理多值函数的情况看来，一旦进入复域，许多隐藏在水下的东西就浮现出来了。所以，数学家一直努力把它涉及的基本思想弄清楚。到 1913 年外尔（H. Weyl）的名著《黎曼曲面的概念》出版，明确指出，黎曼曲面是一个微分流形，而且对于什么是微分流形也给出了今天我们采用的说法。后来的研究表明了，黎曼曲面不能在 $R^3$ 中“实现”，而只有在更高维的 $R^n$ 中才能使它不自交地实现。关于微分流形我们将在第七章中讨论，那时将再回到黎曼曲面的问题。我们现在想要指出的只有一点：函数的定义之内涵不是固定不变的，而是随着数学与物理学的发展，越来越丰富，也越来越深刻地表现了大自然的奥秘。
